\section{CONFIGURAÇÃO EXPERIMENTAL}

\subsection{Ambiente Computacional}
\label{sec:ambiente_computacional}

Os experimentos foram executados em um computador de execução com as seguintes especificações: processador Intel Core i7-14700KF, 32 GB de memória RAM DDR5 a 4533~MT/s (4 módulos Kingston KF552C40-8 de 8~GB cada, formato DIMM), rodando Ubuntu 24.04.03 LTS com Python 3.12.3. O framework de simulação utilizado foi o Flower \cite{flower2020}, em modo de simulação com backend Ray para orquestração dos processos dos clientes federados. Esse modo permite executar servidor e múltiplos clientes em um único processo, preservando o comportamento completo do protocolo federado --- incluindo conversão dos modelos em bytes, transmissão e agregação --- sem exigir infraestrutura de rede real, sendo a abordagem padrão para avaliação experimental de algoritmos de Aprendizado Federado \cite{flower2020}. Os três modelos de \textit{Gradient Boosting} avaliados foram XGBoost \cite{chen2016xgboost}, LightGBM \cite{ke2017lightgbm} e CatBoost \cite{prokhorenkova2018catboost}, e a otimização bayesiana de hiperparâmetros foi realizada com Optuna.

\subsection{Configuração do Dataset}

O dataset BAF \textit{Base} \cite{jesus2022baf} foi preparado conforme descrito na Seção 5, resultando em 999.641 amostras com 99 features (após engenharia de features e \textit{One-Hot Encoding}). A divisão temporal preserva a natureza sequencial das transações: os meses 0--5 formam o conjunto de treinamento (794.680 amostras), o mês 6 o conjunto de validação (108.132 amostras) e o mês 7 o conjunto de teste final (96.829 amostras). A taxa de fraude cresce ligeiramente ao longo do tempo --- 1,03\% no treino, 1,34\% na validação e 1,47\% no teste --- refletindo a evolução natural dos padrões de fraude ao longo dos meses. O severo desbalanceamento de classes no conjunto de treinamento, com 97 transações legítimas para cada fraude, resultou em um \texttt{scale\_pos\_weight} de 96,53, utilizado por todos os algoritmos para compensar o desequilíbrio durante o treinamento. Os dados de treinamento foram distribuídos uniformemente entre 3 clientes, cada um recebendo 264.893 amostras com 2.716 instâncias de fraude (taxa de 1,03\%), em um particionamento IID balanceado que preserva a distribuição original.

\subsection{Otimização Federada de Hiperparâmetros}

Os hiperparâmetros de cada algoritmo foram otimizados de forma federada antes das simulações: cada cliente executou independentemente 50 \textit{trials} de busca bayesiana via Optuna, utilizando 33\% dos seus dados locais para tornar a etapa computacionalmente viável. Os resultados foram agregados no servidor via mediana --- mais robusta a valores atípicos (\textit{outliers}) do que a média aritmética em contextos de dados heterogêneos \cite{miao2022byzantine} --- preservando a localidade dos dados e sem compartilhar amostras entre clientes. A métrica otimizada foi o TPR no ponto de operação de 5\% de FPR, com taxa de aprendizado limitada a 0,05 como \textit{safety brake} para evitar instabilidade no treinamento federado com \textit{warm start}. Os hiperparâmetros resultantes são apresentados na Tabela \ref{tab:hyperparams}.

\begin{table}[H]
\centering
\small
\begin{tabularx}{\textwidth}{|l|l|X|}
\hline
\bf Algoritmo & \bf Hiperparâmetro & \bf Valor \\ \hline
\multirow{6}{*}{XGBoost} & \texttt{max\_depth} & 3 \\ \cline{2-3}
 & \texttt{learning\_rate} & 0,0500 \\ \cline{2-3}
 & \texttt{n\_estimators} & 228 \\ \cline{2-3}
 & \texttt{min\_child\_weight} & 6 \\ \cline{2-3}
 & \texttt{subsample} & 0,7318 \\ \cline{2-3}
 & \texttt{colsample\_bytree} & 0,7858 \\ \hline
\multirow{7}{*}{LightGBM} & \texttt{max\_depth} & 3 \\ \cline{2-3}
 & \texttt{learning\_rate} & 0,0224 \\ \cline{2-3}
 & \texttt{n\_estimators} & 276 \\ \cline{2-3}
 & \texttt{num\_leaves} & 29 \\ \cline{2-3}
 & \texttt{min\_child\_samples} & 54 \\ \cline{2-3}
 & \texttt{subsample} & 0,9310 \\ \cline{2-3}
 & \texttt{colsample\_bytree} & 0,8822 \\ \hline
\multirow{4}{*}{CatBoost} & \texttt{depth} & 4 \\ \cline{2-3}
 & \texttt{learning\_rate} & 0,0500 \\ \cline{2-3}
 & \texttt{iterations} & 298 \\ \cline{2-3}
 & \texttt{l2\_leaf\_reg} & 0,1626 \\ \hline
\end{tabularx}
\caption{Hiperparâmetros otimizados via Optuna federado (agregação por mediana entre 3 clientes)}
\label{tab:hyperparams}
\end{table}

Os três algoritmos convergiram para árvores rasas (\texttt{max\_depth} 3--4), consistente com a literatura de detecção de fraude que indica que modelos mais rasos generalizam melhor em dados com alta proporção de ruído \cite{pozzolo2015calibrating}. A taxa de aprendizado do LightGBM (0,0224) ficou abaixo do limite de 0,05, enquanto XGBoost e CatBoost atingiram exatamente o teto, indicando que o Optuna favoreceu atualizações mais agressivas para esses dois algoritmos.

\subsection{Configuração do Aprendizado Federado}

A Tabela \ref{tab:fl_config} apresenta os parâmetros da simulação federada. A escolha de 50 rodadas de comunicação --- superior às 15--20 rodadas tipicamente reportadas na literatura de FL com boosting \cite{li2020practical,liu2024survey} --- foi deliberada para observar o comportamento de longo prazo de ambas as estratégias e verificar se a convergência sustentada ou a degradação emergiriam com mais rodadas.

\begin{table}[H]
\centering
\begin{tabularx}{\textwidth}{|l|X|}
\hline
\bf Parâmetro & \bf Valor \\ \hline
Número de clientes & 3 \\ \hline
Rodadas de comunicação & 50 \\ \hline
Épocas locais por rodada & 25 \\ \hline
Fração de amostragem & 1,0 (todos os clientes participam) \\ \hline
Particionamento & IID balanceado \\ \hline
Semente aleatória & 42 \\ \hline
FPR alvo & 5\% \\ \hline
\textit{Trials} Optuna por cliente & 50 \\ \hline
\end{tabularx}
\caption{Parâmetros da simulação federada}
\label{tab:fl_config}
\end{table}

Ambas as estratégias utilizam \textit{warm start}: na Bagging, o melhor modelo da rodada anterior é enviado a todos os clientes como ponto de partida; na Cíclica, o modelo percorre os clientes sequencialmente, sendo refinado por cada um antes de passar ao próximo. Os conjuntos de validação e teste são mantidos centralizados no servidor para avaliação global a cada rodada. Ao final do treinamento, o modelo é avaliado no conjunto de teste com dois regimes de limiar: o regime \textbf{standard}, com um único limiar global calibrado para 5\% de FPR no conjunto completo, e o regime \textbf{fair}, com limiares independentes por grupo demográfico (idade $\geq 50$ e idade $< 50$), cada um calibrado para 5\% de FPR no respectivo grupo --- técnica de \textit{post-processing} que equaliza as taxas de falso positivo sem modificar o modelo treinado \cite{hardt2016equality}.

\subsection{Experimentos Realizados}

Foram conduzidas 6 simulações, cobrindo as combinações de 3 algoritmos com 2 estratégias de agregação (Bagging e Cíclica). O experimento completo --- incluindo pré-processamento, otimização de hiperparâmetros e as 6 simulações --- foi executado em 6.251 segundos (aproximadamente 1 hora e 44 minutos) no computador de execução descrito na Seção~\ref{sec:ambiente_computacional}. As métricas reportadas para cada configuração são: TPR @ 5\% FPR no conjunto de teste, AUC-ROC, FPR Ratio ($\min/\max$ das FPR entre grupos demográficos, onde 1,0 indica fairness perfeita), tempo de execução da simulação e volume total de dados transferidos entre clientes e servidor.
